% ----------------------------------------------------------------------
% Parte II, capítulo 5
% ----------------------------------------------------------------------
\chapter{Resultados e discussão}
\label{cap:resultados}

\instrucao{Apresente os resultados por dimensão do modelo e discuta cada um à luz da Parte~\ref{part:fundamentos}. Separe o que os dados mostram do que você interpreta.}

\section{Índices por dimensão}
A Tabela~\ref{tab:indices} resume os índices obtidos (dados de exemplo), calculados pela Equação~\eqref{eq:indice}.

\begin{table}[htb]
  \centering
  \caption{Índices de interoperabilidade por região}
  \label{tab:indices}
  \begin{tabular}{lcccc}
    \toprule
    \multirow{2}{*}{\textsc{Região}} & \multicolumn{3}{c}{\textsc{Dimensão}} & \multirow{2}{*}{$I$} \\
    \cmidrule(lr){2-4}
     & $I_s$ & $I_m$ & $I_o$ & \\
    \midrule
    Norte & 0,91 & 0,42 & 0,55 & 0,63 \\
    Nordeste & 0,94 & 0,47 & 0,61 & 0,67 \\
    Sudeste & 0,97 & 0,58 & 0,70 & 0,75 \\
    \bottomrule
  \end{tabular}
  \fonte{Dados da pesquisa (exemplo). Pesos: $w_s = 0{,}3$; $w_m = 0{,}4$; $w_o = 0{,}3$.}
\end{table}

\section{Discussão}
\instrucao{Volte à tese defendida na introdução: os dados a sustentam, a refutam ou pedem ajuste?}
Em todas as regiões, a dimensão sintática supera 0,9, enquanto a semântica fica abaixo de 0,6. O resultado sustenta a tese de que o gargalo está no controle de autoridades e no uso de vocabulários, e não no protocolo de troca.
